% !TEX root = ../main.tex
% Conclusion — drafted per manuscript/planning/conclusion_drafting_spec.md
% Composition: 1α + 2α + 3α + 4β + Fusion A (5E.1 + 5F.2). ≈ 380 w.
% No equations, no figures, no \cref to figures. Two \cite to keys already
% present in references.bib and already cited in results.tex / discussion.tex.

\section{Conclusions}%
\label{Conclusions}

The extensions of general relativity that gravitational-wave
observatories will probe introduce many new interactions, whose effect
on the Gertsenshtein channel has not, until now, been systematically
explored. We have presented \TIDAL{}, a symbolic--numerical software
package that derives the linearised equations of motion directly from
the Lagrangian, and applied it to the quadratic Poincar\'{e}-gauge
action coupled to electromagnetism.

Across the surveyed operator space we find that the Gertsenshtein
conversion responds to the full coupling content of each Lagrangian,
but a few operators are robustly dominant --- amplifying the
conversion no matter what other terms are present. The most consistent
of these is the derivative-mixing contraction
\begin{equation}
    \lag \supset \rcD{^\alpha}\MAGT{^{\beta\gamma\delta}}\, \MAGF{_{\alpha\beta\gamma\delta}},
\end{equation}
which is also the dominant operator-level distinction between
amplifying and suppressing parameter regions; for most other operators
those regions largely overlap. The amplification persists without
propagating-torsion, so purely ghost-driven growth cannot explain
it.\wb{That is a nice result}

Whether the exponential mode growth behind these amplifications
corresponds to a physical amplification, a tachyonic instability, or a
wrong-sign kinetic ghost is a question the present survey does not
address. A particle-spectrum analysis via
\PSALTer{}~\cite{barker2024psalter,barker2025psalter2} would determine
which of the three applies for each amplifying direction.

Re-running the survey on a finite-extent magnetic background would
give the propagating wave only a finite interaction time, bounding the
exponential growth that otherwise compounds indefinitely. The sectors
that overflow under that embedding would then re-enter the survey at
finite amplitude, and the present coupling-level identifications would
be replaced by physical conversion magnitudes comparable with current
indirect-detection thresholds for the Gertsenshtein signal. A second
extension is to sample the direction of each coupling vector
independently of its magnitude (\cref{CubedSphereAtlas}); this would
identify which combinations of operators drive the amplification,
independently of the basis in which the action is written.

Gertsenshtein channel amplification is distributed across the operator
content of the action beyond Einstein--Maxwell, with a few robustly
dominant operators carrying the conversion no matter what else is
present. Because the amplification persists without propagating
torsion, future analyses can sidestep the difficult problem of which
modes propagate and focus directly on the operator content of the
action. Whether each amplifying coupling identified here corresponds
to a physical amplification or to a vacuum-dependent instability is a
question the survey cannot, by construction, decide; what it does
decide is which operators of the quadratic gauge-gravity action are
candidates for that question at all --- and therefore which couplings
any future bound on the Gertsenshtein signal would target.

\wb{OK, so the main thing that is missing here is a concrete discussion of how large the effects are in real terms. You open with astrophysics, and with scales and numbers, but most of the results are expressed somewhat obscurely in terms of evidence. Is there a concrete connection between the plots and the degree to which the effect has been amplified? How can an astrophysicist decide whether a model selected from your catalogue will provide effects at a specified physical scale? You probably want to add a `translation' paragraph here in the conclusions, that allows someone to take a point from the plots and say `this is the degree of amplification I get for a magnetic field of such-and-such strength, and a gravitational wave of such-and-such frequency'.}